\subsection{Two-Sample Normal Inference}

Two-sample inference compares the location of two response distributions.  In a controlled experiment, one sample may receive a treatment and the other a control; in an observational comparison, the same formulas may be computed but their causal interpretation depends on the design.

\begin{definition} \label{def:two-independent-sample-location-model} \hlt{Two Independent Samples and the Location Contrast}\\
$X_1,\ldots,X_n\sim F$ are independent, $Y_1,\ldots,Y_m\sim G$ are independent, and the two samples are independent.
\begin{enumerate}[label=(\roman*)]
\setlength{\itemsep}{0pt}
\item \hlt{Means}: $\mu_X\defeq\mathbb{E}X_i$ and $\mu_Y\defeq\mathbb{E}Y_j$ summarize the two average levels, assuming both exist.
\item \hlt{Location contrast}: $\Delta\defeq\mu_X-\mu_Y$ measures the shift from group $Y$ to group $X$.
\item \hlt{Estimator}: $\widehat{\Delta}\defeq\overline{X}-\overline{Y}$ estimates $\Delta$, where $\overline{X}\defeq n^{-1}\sum\limits_iX_i$ and $\overline{Y}\defeq m^{-1}\sum\limits_jY_j$.
\end{enumerate}
\end{definition}

\begin{proposition} \label{prop:two-sample-known-variance-normal-pivot} \hlt{Known Common Variance Pivot}\\
Let $X_1,\ldots,X_n$ and $Y_1,\ldots,Y_m$ be independent samples, $X_i\sim N(\mu_X,\sigma^2)$, $Y_j\sim N(\mu_Y,\sigma^2)$, known $\sigma>0$.\\
Set $\Delta\defeq\mu_X-\mu_Y$, $\widehat{\Delta}\defeq\overline{X}-\overline{Y}$, and $\operatorname{se}_0\defeq\sigma\sqrt{n^{-1}+m^{-1}}$.
\begin{enumerate}[label=(\roman*)]
\setlength{\itemsep}{0pt}
\item $\widehat{\Delta}\sim N(\Delta,\operatorname{se}_0^2)$ and $Z\defeq\frac{\widehat{\Delta}-\Delta}{\operatorname{se}_0}\sim N(0,1)$.
\item $C_\alpha(X,Y)\defeq[\widehat{\Delta}-z_{1-\frac{\alpha}{2}}\operatorname{se}_0,\widehat{\Delta}+z_{1-\frac{\alpha}{2}}\operatorname{se}_0]$ satisfies $\mathbb{P}_\Delta\{\Delta\in C_\alpha(X,Y)\}=1-\alpha$.
\end{enumerate}
\begin{proof}
\begin{enumerate}[label=(\roman*)]
\setlength{\itemsep}{0pt}
\item Since $\overline{X}\sim N(\mu_X,\frac{\sigma^2}{n})$, $\overline{Y}\sim N(\mu_Y,\frac{\sigma^2}{m})$, and $\overline{X}\perp\!\!\!\perp\overline{Y}$, $\widehat{\Delta}=\overline{X}-\overline{Y}\sim N(\Delta,\frac{\sigma^2}{n}+\frac{\sigma^2}{m})=N(\Delta,\operatorname{se}_0^2)$.\
Standardization gives $Z\sim N(0,1)$.
\item
By (i), $\mathbb{P}_\Delta\{\abs{Z}\leq z_{1-\frac{\alpha}{2}}\}=1-\alpha$, equivalently $\mathbb{P}_\Delta\{\Delta\in C_\alpha(X,Y)\}=1-\alpha$.
\end{enumerate}
\hlt{Structural Commentary}: The contrast is a linear combination of independent sample means.\\
\hlt{Key Insight}: Known common variance gives an exact normal pivot for $\Delta$.
\end{proof}
\end{proposition}

\begin{definition} \label{def:pooled-two-sample-variance} \hlt{Pooled Sample Variance}\\
Let $X_1,\ldots,X_n\sim N(\mu_X,\sigma^2)$ and $Y_1,\ldots,Y_m\sim N(\mu_Y,\sigma^2)$ be independent samples with common unknown $\sigma^2$.
\begin{enumerate}[label=(\roman*)]
\setlength{\itemsep}{0pt}
\item \hlt{Sample variances}: $S_X^2\defeq(n-1)^{-1}\sum\limits_{i=1}^n(X_i-\overline{X})^2$ and $S_Y^2\defeq(m-1)^{-1}\sum\limits_{j=1}^m(Y_j-\overline{Y})^2$.
\item \hlt{Pooled variance}: $S_p^2\defeq\frac{(n-1)S_X^2+(m-1)S_Y^2}{n+m-2}$.
\item \hlt{Degrees-of-freedom weighting}: $S_p^2$ weights $S_X^2$ and $S_Y^2$ by $n-1$ and $m-1$, respectively.
\end{enumerate}
\end{definition}

\begin{theorem} \label{thm:two-sample-pooled-t-pivot} \hlt{Pooled Two-Sample Student Pivot}\\
Let $X_i\sim N(\mu_X,\sigma^2)$, $i=1,\ldots,n$, and $Y_j\sim N(\mu_Y,\sigma^2)$, $j=1,\ldots,m$, be mutually independent, with unknown common $\sigma^2>0$. Set $\Delta\defeq\mu_X-\mu_Y$ and $\nu\defeq n+m-2$. Then $T\defeq\frac{(\overline{X}-\overline{Y})-\Delta}{S_p\sqrt{n^{-1}+m^{-1}}}\sim t_\nu$.
\begin{proof}
Let $Z\defeq\frac{(\overline{X}-\overline{Y})-\Delta}{\sigma\sqrt{n^{-1}+m^{-1}}}$ and $W\defeq\frac{\nu S_p^2}{\sigma^2}=\frac{(n-1)S_X^2+(m-1)S_Y^2}{\sigma^2}$.\\
Since $\overline{X}-\overline{Y}\sim N(\Delta,\sigma^2(n^{-1}+m^{-1}))$, $Z\sim N(0,1)$. Let $U_X\defeq\frac{(n-1)S_X^2}{\sigma^2}$ and $U_Y\defeq\frac{(m-1)S_Y^2}{\sigma^2}$.\\
By normal mean-residual independence and sample independence, $Z,U_X,U_Y$ mutually independent,\\
with $U_X\sim\chi^2_{n-1}$, $U_Y\sim\chi^2_{m-1}$. Hence $W=U_X+U_Y\sim\chi^2_\nu$ and $W\perp\!\!\!\perp Z$. Since $S_p^2=\frac{\sigma^2W}{\nu}$, $T=\frac{Z}{\sqrt{\frac{W}{\nu}}}\sim t_\nu$.\\
\hlt{Structural Commentary}: The pooled variance combines the two residual sums of squares.\\
\hlt{Key Insight}: A common variance permits replacing $\sigma$ by an independent pooled estimate.
\end{proof}
\end{theorem}

\begin{method} \label{meth:pooled-two-sample-t-interval-and-tests} \hlt{Pooled Two-Sample $t$ Procedures}\\
Under the equal-variance normal model, set $\Delta\defeq\mu_X-\mu_Y$, $\nu\defeq n+m-2$, and $\operatorname{se}_p\defeq S_p\sqrt{n^{-1}+m^{-1}}$.
\begin{enumerate}[label=(\roman*)]
\setlength{\itemsep}{0pt}
\item \hlt{Interval}: $C_\alpha\defeq[(\overline{X}-\overline{Y})-t_{\nu,1-\frac{\alpha}{2}}\operatorname{se}_p,(\overline{X}-\overline{Y})+t_{\nu,1-\frac{\alpha}{2}}\operatorname{se}_p]$ is exact $1-\alpha$ for $\Delta$.
\item \hlt{Statistic}: for $H_0:\Delta=0$, use $t_0\defeq\frac{\overline{X}-\overline{Y}}{\operatorname{se}_p}$, with $t_0\sim t_\nu$ under $H_0$.
\item \hlt{Rejection}: reject $H_0$ for $\abs{t_0}>t_{\nu,1-\frac{\alpha}{2}}$ under $H_1:\Delta\neq0$\\
reject $H_0$ for $t_0>t_{\nu,1-\alpha}$ under $H_1:\Delta>0$\\
reject $H_0$ for $t_0<-t_{\nu,1-\alpha}$ under $H_1:\Delta<0$.
\item \hlt{Duality}: the two-sided test rejects iff $0\notin C_\alpha$.
\end{enumerate}
\end{method}

\begin{example} \label{ex:latent-heat-two-methods} \hlt{Latent Heat of Fusion: Two Methods}\\
For methods $A$ and $B$, let $n_A=13$, $n_B=8$, $\overline{x}_A=80.02$, $\overline{x}_B=79.98$, $s_A=0.024$, $s_B=0.031$, and $S_p=0.027$.\\
Then $\widehat{\Delta}=\overline{x}_A-\overline{x}_B=0.04$, $\widehat{\operatorname{se}}=S_p\sqrt{13^{-1}+8^{-1}}\approx0.012$, pooled $95\%$ interval for $\Delta=\mu_A-\mu_B$ is\\
$\approx (0.015,0.065)$. For $H_0:\Delta=0$, $t_0=\frac{\widehat{\Delta}}{\widehat{\operatorname{se}}}\approx3.33$ with $\nu=19$, strong evidence that method means differ.
\end{example}

\begin{proposition} \label{prop:two-sample-normal-glrt-equivalent-to-t} \hlt{Likelihood-Ratio View of the Pooled Test}\\
Let $X_i\sim N(\mu_X,\sigma^2)$, $i=1,\ldots,n$, and $Y_j\sim N(\mu_Y,\sigma^2)$, $j=1,\ldots,m$, be mutually independent, with unknown common $\sigma^2>0$.
Set $\Delta\defeq\mu_X-\mu_Y$, $\nu\defeq n+m-2$, and $T\defeq\frac{\overline{X}-\overline{Y}}{S_p\sqrt{n^{-1}+m^{-1}}}$.\\
For $H_0:\Delta=0$ against $H_1:\Delta\neq0$, $\lambda(X,Y)=(1+\frac{T^2}{\nu})^{-\frac{n+m}{2}}$.\\
Hence the GLR rejects for large $\abs{T}$; at level $\alpha$, reject iff $\abs{T}>t_{\nu,1-\frac{\alpha}{2}}$.
\begin{proof} Set $N\defeq n+m$, $\nu\defeq N-2$, $d\defeq\overline{X}-\overline{Y}$, $R_1\defeq\sum\limits_i(X_i-\overline{X})^2+\sum\limits_j(Y_j-\overline{Y})^2$, $\widehat{\mu}_0\defeq\frac{n\overline{X}+m\overline{Y}}{N}$, and $R_0\defeq\sum\limits_i(X_i-\widehat{\mu}_0)^2+\sum\limits_j(Y_j-\widehat{\mu}_0)^2$. Under $H_1$, $(\widehat{\mu}_X,\widehat{\mu}_Y,\widehat{\sigma}_1^2)=(\overline{X},\overline{Y},\frac{R_1}{N})$; under $H_0$, $(\widehat{\mu},\widehat{\sigma}_0^2)=(\widehat{\mu}_0,\frac{R_0}{N})$.\\
The maximised Gaussian likelihood is proportional to $\widehat{\sigma}^{-N}$, hence $\lambda=(\frac{R_1}{R_0})^{\frac{N}{2}}$.\\
ANOVA decomposition gives $R_0=R_1+\{\frac{nm}{N}\}d^2$, while $T^2=\frac{d^2}{(\frac{R_1}{\nu})(n^{-1}+m^{-1})}=\frac{\nu\{\frac{nm}{N}\}d^2}{R_1}$.\\
Therefore $\frac{R_0}{R_1}=1+\frac{T^2}{\nu}$ and $\lambda=(1+\frac{T^2}{\nu})^{-\frac{N}{2}}$, decreasing in $T^2$; small $\lambda$ iff large $\abs{T}$.\\
\hlt{Structural Commentary}: The null fit adds only the between-sample mean discrepancy to the unrestricted fit.\\
\hlt{Key Insight}: The GLR depends on the data through the pooled $t$ statistic.
\end{proof}
\end{proposition}

\begin{method} \label{meth:welch-two-sample-test} \hlt{Unequal-Variance Two-Sample Procedure}\\
For independent samples without the equal-variance assumption, set $\widehat{v}\defeq \frac{S_X^2}{n}+\frac{S_Y^2}{m}$.
\begin{enumerate}[label=(\roman*)]
\setlength{\itemsep}{0pt}
\item \hlt{Standard error}: estimate $\operatorname{Var}(\overline{X}-\overline{Y})$ by $\widehat{v}$.
\item \hlt{Welch statistic}: for $H_0:\Delta=\Delta_0$, use $T_W\defeq\frac{(\overline{X}-\overline{Y})-\Delta_0}{\sqrt{\widehat{v}}}$.
\item \hlt{Degrees of freedom}: calibrate $T_W$ by $t_\nu$, where $\nu\defeq\frac{\widehat{v}^2}{\frac{(S_X^2/n)^2}{n-1}+\frac{(S_Y^2/m)^2}{m-1}}$.
\item \hlt{Use}: the procedure avoids pooling variances, at the cost of approximate calibration.
\end{enumerate}
\end{method}

\begin{corollary} \label{cor:two-normal-sample-variance-ratio-f} \hlt{Variance Ratio Test Under Normality}\\
Let $X_1,\ldots,X_n\sim N(\mu_X,\sigma_X^2)$ and $Y_1,\ldots,Y_m\sim N(\mu_Y,\sigma_Y^2)$ be independent samples, with $n,m\geq2$.\\
Set $\rho\defeq\frac{\sigma_X^2}{\sigma_Y^2}$ and $R\defeq \frac{S_X^2}{S_Y^2}$. Then $\frac{R}{\rho}\sim F_{n-1,m-1}$.
\begin{proof} Let $U\defeq\frac{(n-1)S_X^2}{\sigma_X^2}$ and $V\defeq\frac{(m-1)S_Y^2}{\sigma_Y^2}$.\\
Normal residual chi-square laws give $U\sim\chi^2_{n-1}$ and $V\sim\chi^2_{m-1}$, and sample independence gives $U\perp\!\!\!\perp V$.\\
Thus $\frac{\frac{U}{n-1}}{\frac{V}{m-1}}=\frac{R}{\rho}\sim F_{n-1,m-1}$. If $H_0:\rho=1$, then $R\sim F_{n-1,m-1}$, yielding exact normal-theory tests and confidence intervals for $\rho$; the calibration is normality-sensitive.\\
\hlt{Structural Commentary}: The comparison standardises two independent residual variance estimates.\\
\hlt{Key Insight}: Normality turns each sample variance into an independent chi-square pivot.
\end{proof}
\end{corollary}

\begin{method} \label{meth:two-sample-normal-power-balanced} \hlt{Approximate Power for Balanced Two-Sample $z$ Tests}\\
For planning, assume balanced samples of size $n$, known common $\sigma$, and target effect $\Delta\defeq\mu_X-\mu_Y>0$.
\begin{enumerate}[label=(\roman*)]
\setlength{\itemsep}{0pt}
\item \hlt{Statistic}: use $Z\defeq\frac{\overline{X}-\overline{Y}}{\sigma\sqrt{\frac{2}{n}}}$ and reject two-sided at level $\alpha$ when $\abs{Z}>z_{1-\frac{\alpha}{2}}$.
\item \hlt{Shift}: under $\Delta>0$, $Z\sim N(\delta,1)$ with $\delta\defeq(\frac{\Delta}{\sigma})\sqrt{\frac{n}{2}}$.
\item \hlt{Power}: $\pi(\Delta)\defeq1-\Phi(z_{1-\frac{\alpha}{2}}-\delta)+\Phi(-z_{1-\frac{\alpha}{2}}-\delta)$.
\item \hlt{Sample size}: if the lower-tail term is negligible, target power $1-\beta$ gives $n\approx\frac{2\sigma^2(z_{1-\frac{\alpha}{2}}+z_{1-\beta})^2}{\Delta^2}$.
\end{enumerate}
\end{method}

\begin{figure}[H]
  \centering
  \begin{tikzpicture}[
  x=0.72cm,
  y=0.72cm,
  >=Stealth,
  line cap=round,
  line join=round,
  every node/.style={font=\scriptsize},
  axis/.style={->, gray!70, line width=0.45pt},
  tick/.style={gray!70, line width=0.4pt},
  curve/.style={blue!70!black, line width=0.95pt},
  guide/.style={gray!45, dashed, line width=0.45pt},
  label/.style={text=black!78}
]
  \draw[axis] (0,0) -- (6.45,0) node[black!78,right] {$\Delta$};
  \draw[axis] (0,0) -- (0,3.45) node[black!78,above] {power};
  \foreach \y/\lab in {0/0,0.6/.2,1.8/.6,3/1.0} {
    \draw[tick] (-0.08,\y) -- (0.08,\y) node[black!78,left=5pt] {$\lab$};
  }
  \foreach \x/\lab in {0/0,2/2,4/4,6/6} {
    \draw[tick] (\x,-0.08) -- (\x,0.08) node[black!78,below=5pt] {$\lab$};
  }
  \draw[curve, smooth]
    plot coordinates {
      (0.0,0.14) (0.5,0.24) (1.0,0.45) (1.5,0.80)
      (2.0,1.22) (2.5,1.68) (3.0,2.10) (3.5,2.45)
      (4.0,2.68) (4.5,2.83) (5.0,2.92) (5.5,2.97)
      (6.0,2.99)
    };
  \draw[guide] (0,2.7) -- (6.1,2.7);
  \node[label, anchor=west] at (4.25,2.48) {$0.9$};
\end{tikzpicture}

  \caption{Power rises as the detectable mean difference $\Delta$ increases.}
  \label{fig:two-sample-power-curve}
\end{figure}

\begin{remark} \label{rem:two-sample-transformations-and-diagnostics} \hlt{Transformations and Normal Diagnostics}\\
Normal-theory two-sample procedures are often CLT-robust for large samples, but small skewed or outlier-prone samples require diagnostics. Transformations may stabilize spread and improve symmetry, but they change the estimand, e.g. a log-scale mean difference corresponds to a ratio-scale comparison.
\end{remark}
